\chapter{Lexical structure}
\label{ch:lexical}

A Forge source file is a sequence of ASCII characters, read as tokens
separated by whitespace (spaces, tabs, line feeds and carriage returns) and
comments.

\section{Comments}
\index{comments}
A comment is either a line comment, from \fc|//| to the end of the line, or
a block comment between \fc|/*| and \fc|*/|. Block comments do not nest.

\section{Identifiers and keywords}
\index{identifiers}
An identifier starts with a letter or an underscore, followed by letters,
digits and underscores: \fc|x|, \fc|list_length|, \fc|_tmp2|. Identifiers that
start with an upper-case letter name constructors (\fc|Some|, \fc|Node|),
types (\fc|Option|) and modules (\fc|Geometry|). A qualified name
\fc|M.x| refers to \fc|x| inside module \fc|M| (Section~\ref{sec:modules}).
Identifiers beginning with two underscores are used internally and should be
avoided.

The following words are reserved (Appendix~\ref{ch:keywords} explains each):

\begin{center}
\small\ttfamily
\begin{tabular}{@{}llllllll@{}}
and & case & do & else & end & enum & false & fn \\
for & if & impl & import & in & inter & let & macro\_rules \\
mod & of & rec & requires & switch & then & trait & true \\
type & use & val & where & with & & & \\
\end{tabular}
\end{center}

\noindent The type names \fc|Int|, \fc|Float|, \fc|Bool|, \fc|Char|,
\fc|String| and \fc|Unit| are also keywords. The lower-case spellings
\texttt{int}, \texttt{float}, \texttt{bool}, \texttt{char}, \texttt{string}
and \texttt{unit} are accepted in types as synonyms.

\section{Literals}
\index{literals}

\begin{description}[style=nextline,leftmargin=2em,font=\normalfont\sffamily\bfseries]
  \item[Integers] A sequence of decimal digits: \fc|0|, \fc|42|. A literal must fit
    in a machine integer (\S\ref{sec:int-width}); \fc|-17| is the negation of
    \fc|17|.
  \item[Floats] Digits, a dot and optional digits: \fc|3.14|, \fc|5.|. There is no
    exponent notation.
  \item[Characters] A single character or an escape between single quotes:
    \fc|'a'|, \fc|'\n'|, \fc|'\''|.
  \item[Strings] Characters and escapes between double quotes:
    \fc|"line\n"|. A line break inside the quotes is part of the string.
  \item[Booleans and unit] \fc|true|, \fc|false| and \fc|()|.
\end{description}

The escapes are \verb|\n| (line feed), \verb|\t| (tab), \verb|\r| (carriage
return), \verb|\0| (NUL), \verb|\\|, \verb|\"| and \verb|\'|; any other
backslash sequence is an error.

\forgeexample{lex-literals}

\section{Operators and punctuation}
\index{operators}
The fixed operators and punctuation are
\fc@+ - * / % ^ :: == != <> < > <= >= && || |> -> <- = | , ; : . ...@
and the brackets \texttt{( ) [ ] \{ \}}.
Any other sequence of the characters \fc@+ - * / % < > = ^ ! |@ that
starts with one of \fc@+ - * / % < > = ^ !@ is a \emph{custom operator},
which a program defines like a function, \fc|let (+++) a b = ...|; its
precedence is that of its first character (Section~\ref{sec:operators}).
The characters \fc|&| and \fc!|! begin only the fixed operators.
